\section{Recursos do abntexto-ufc, elementos pós-textuais e revisão final}\label{sec:recursos-finais}

A classe reúne opções de perfil e módulos editoriais para que o mesmo projeto possa atender diferentes tipos de trabalho sem duplicar a lógica de formatação. Esta seção encerra o guia explicando os principais recursos e os elementos que aparecem depois do texto principal.

\subsection{Perfis de documento}

A opção \texttt{type} seleciona o perfil acadêmico. A matriz de certificação atual exercita \texttt{undergraduate-capstone}, \texttt{specialization-capstone}, \texttt{masters-thesis}, \texttt{doctoral-thesis}, \texttt{research-project} e \texttt{anonymized-research-project}.

\guiapolitica{Perfis controlam a aplicabilidade de elementos e metadados. O objetivo é evitar que o usuário tenha de remover manualmente partes da classe para adaptar um TCC a uma dissertação, tese ou projeto.}

\guianormativa{ABNT NBR 15287:2025}{Os perfis \texttt{research-project} e \texttt{anonymized-research-project} usam a estrutura de projeto de pesquisa adotada pelo projeto. O perfil anonimizado adiciona uma política técnica de ocultação de dados pessoais para processos que exijam avaliação sem identificação; essa anonimização não é apresentada como regra geral da NBR 15287 \parencite{abnt15287}.}

\subsection{Fontes e motores de compilação}

A opção \texttt{font=times|arial} seleciona a família textual desejada. A opção \texttt{strict-font} determina se a identidade literal é obrigatória ou se o ambiente pode usar fallback portátil.

\guiainstitucional{A política tipográfica UFC adotada pelo projeto trabalha com Arial ou Times New Roman. A certificação literal é executada em Windows com pdfLaTeX e LuaLaTeX; as fontes Microsoft não são redistribuídas pelo repositório.}

\guiapolitica{O documento de referência usa modo portátil para compilar em ambientes públicos. Uma compilação com fallback não deve ser confundida com evidência de identidade literal da família solicitada.}

\subsection{PDF/A}

O cabeçalho de \texttt{main.tex} declara \texttt{pdfstandard=A-2b}. A suíte verifica o PDF gerado com veraPDF e também confere incorporação de fontes.

\guiainstitucional{O fluxo de depósito acompanhado pelo projeto exige PDF/A nos casos institucionais aplicáveis. As páginas atuais do Sistema de Bibliotecas devem ser verificadas conforme o tipo de trabalho \parencite{ufcdepositotcc2026,ufcdepositopos2026}.}

\guiapolitica{A escolha do perfil específico PDF/A-2b para o documento de referência e para a certificação é uma política técnica do projeto; ela não é apresentada como uma afirmação de que toda norma ABNT exige especificamente PDF/A-2b.}

\subsection{Referências}

As referências aparecem imediatamente após o texto quando produzidas por \texttt{\string\ufcPrintReferences}. O banco bibliográfico é declarado com \texttt{\string\ufcAddBibliographyResource} e processado por Biber no fluxo recomendado.

\guianormativa{ABNT NBR 6023:2025}{Inclua na lista final as fontes de acordo com a política bibliográfica do trabalho e mantenha os dados necessários à identificação de cada documento \parencite{abnt6023}.}

\subsection{Glossário}

O glossário é um elemento opcional para definir termos especializados, siglas ou conceitos quando essa lista realmente auxiliar o leitor. O documento de referência o mantém habilitado para demonstrar a integração com o pacote \texttt{glossaries}.

\guiapolitica{Não use o glossário como duplicação automática de toda palavra técnica. Registre apenas termos cuja definição agregue valor ao trabalho e cuja repetição no corpo prejudicaria a leitura.}

\subsection{Apêndices}

Apêndices contêm material complementar elaborado pelo próprio autor e são identificados em sequência própria.

\guianormativa{estrutura pós-textual adotada pelo projeto}{Cada apêndice inicia em página própria. No escopo auditado, seu cabeçalho é centralizado, em letras maiúsculas, negrito e tamanho 12.}

\guiaexemplo{Este documento possui apêndices com material textual, questionário, código-fonte e orientação para documento externo autoral. Eles existem para exercitar diferentes comprimentos e conteúdos.}

\subsection{Anexos}

Anexos são usados para material complementar de autoria externa ou não produzido pelo autor do trabalho.

\guianormativa{estrutura pós-textual adotada pelo projeto}{Cada anexo inicia em página própria e usa a mesma política de apresentação do cabeçalho exercitada para apêndices.}

\guiaexemplo{Se um documento externo em PDF precisar ser anexado, verifique legibilidade, direitos de uso, orientação de página e compatibilidade com o arquivo final. O exemplo deste projeto documenta a rota sem redistribuir material externo desnecessário.}

\subsection{Índice}

O índice remissivo é diferente do sumário: ele organiza termos e assuntos para facilitar a localização de ocorrências no documento.

\guianormativa{ABNT NBR 6034:2004}{O índice é elemento opcional e possui norma específica de apresentação. A classe demonstra o recurso com \texttt{imakeidx} e verifica o título e sua composição no PDF final \parencite{abnt6034}.}

\subsection{Lombada}

\guianormativa{ABNT NBR 12225:2023}{A lombada possui aplicação condicional, normalmente relacionada à produção física do trabalho. O documento eletrônico de referência não a força em perfis nos quais não seja aplicável \parencite{abnt12225}.}

\subsection{Checklist antes da entrega}

Antes de entregar um trabalho real produzido a partir deste modelo:

\begin{enumerate}
  \item revise todos os campos de \texttt{\string\ufcsetup};
  \item mantenha apenas os pré-textuais aplicáveis ao seu perfil;
  \item substitua todos os exemplos e conteúdos de demonstração pelo texto acadêmico real;
  \item confirme citações e referências contra as fontes consultadas;
  \item revise figuras, tabelas, quadros, códigos, algoritmos, equações, fontes e notas;
  \item verifique glossário, apêndices, anexos e índice somente quando pertinentes;
  \item compile o documento completo com todos os processadores auxiliares necessários;
  \item examine visualmente o PDF final, inclusive quebras de página e elementos externos;
  \item execute os gates de validação disponíveis no projeto;
  \item confira as regras atuais do curso, programa e fluxo de depósito da UFC.
\end{enumerate}

\guiapolitica{Um CI verde demonstra que o corpus testado respeitou os contratos executáveis do projeto. Ele não transforma automaticamente toda regra externa em prova formal de conformidade e não substitui a revisão acadêmica e institucional do conteúdo.}

Este modelo comentado deve ser entendido como ponto de partida auditável: o usuário escreve o conteúdo acadêmico, enquanto a classe concentra a maior parte das decisões editoriais e o repositório registra a evidência usada para mantê-las.

\index{PDF/A}\index{glossário}\index{apêndice}\index{anexo}\index{índice}\index{projeto de pesquisa}
